% =========================================================
% CSC311 Final Project Report Template
% Single column, 12pt, 8 pages max.
% =========================================================
\documentclass[12pt]{article}
% ---------- Packages ----------
\usepackage[margin=1in]{geometry}
\usepackage{parskip} % blank line between paragraphs, no indent
\usepackage{graphicx}
\usepackage{booktabs}
\usepackage{amsmath, amssymb}
\usepackage{siunitx}
\usepackage[hidelinks]{hyperref}
\usepackage{caption}
\usepackage{enumitem}
\usepackage{xcolor}
\usepackage{XCharter}
\usepackage{float}
\usepackage[xcharter]{newtxmath}
\usepackage[scaled=0.96]{inconsolata}
%--------- Headers and footers ------
\usepackage{fancyhdr}
\pagestyle{fancy}
\setlength{\headheight}{14.49998pt}
\fancyhf{}
\fancyhead[L]{CSC311 Project Report}
\fancyhead[R]{Group Members}
\fancyfoot[C]{\thepage}
% ---------- Captions & Lists ----------
\captionsetup{font=small, labelfont=bf}
\setlist[itemize]{noitemsep, topsep=2pt}
\setlist[enumerate]{noitemsep, topsep=2pt}
% ---------- Helpful macros ----------
\newcommand{\bestmodel}{\textit{[Best Model Name]}}
\newcommand{\projacc}{\textbf{[XX.XX\%]}}
% ---------- Safe figure helper ----------
\makeatletter
\newcommand{\maybeincludegraphics}[3]{%
\begin{figure}[h]
\centering
\IfFileExists{#1}{%
\includegraphics[width=0.7\linewidth]{#1}%
}{%
\fbox{%
\parbox[c][0.28\textheight][c]{0.7\linewidth}{%
\centering
\vspace{0.5em}
\textit{[Placeholder for \texttt{#1}]}\\
Add your figure here.
\vspace{0.5em}
}%
}%
}%
\caption{#2}
\label{#3}
\end{figure}%
}
\makeatother
\usepackage{xcolor}
\definecolor{royalblue}{rgb}{0.0, 0.14, 0.4}
\definecolor{royalblue(web)}{rgb}{0.25, 0.41, 0.88}
\usepackage{tcolorbox}
\tcbuselibrary{breakable}
\usepackage{comment}
\usepackage{tikz}
\usetikzlibrary{arrows.meta}
% =========================================================
% Title Page
% =========================================================
\begin{document}
\begin{titlepage}
\centering
\vspace*{3cm}
{\Large {CSC311: Introduction to Machine Learning} \par}
\vspace{1cm}
{\Large {Final Project Report} \par}
\vspace{1cm}
% {\large \textit{Optional Creative Project Title} \par}
\vspace{3cm}
{\large Submitted by: Bianca Bian, Sylvia Wong, Chris Yang \par}
{\large Group Number: 63654 \par}
\vspace{0.5cm}
\vfill
{\large University of Toronto \\ Fall 2026 \par}
\end{titlepage}
\clearpage
\setcounter{page}{1}
% ================================================
% 2. Executive Summary
% ================================================
\section{Executive Summary}
% This section should be at most 1/4 page (1 paragraph).
% \begin{itemize}
% \item Summarize the model families you explored (at least three).
% \item State your best-performing model and its projected accuracy.
% \item Provide a high-level justification of why this model outperformed the
% others.
% \end{itemize}

The following report describes the data exploration process, feature processing, model families investigated, and results of the best-performing model to classify sports based on survey responses. During the model development, we compared the performance of softmax regression, random forests, an ensemble between Naive Bayes and random forests on the validation set. The best-performing model was a random forest with a depth of 15, achieving a validation accuracy of 0.95. Random forest outperformed softmax regression and Naive Bayes due to its ability to capture non-linear relationships and interactions between features. Softmax regression failed to capture complex interactions between features and the poor performance of Naive Bayes resulted in a low accuracy in the ensemble model. We expect the Random Forest model to achieve a projected accuracy of 0.90 on the unseen test set.

\section{Data Exploration}\label{sec:data-exploration}

\subsection*{Dataset Overview}
\begin{itemize}[leftmargin=1.5em]
    \item \textbf{Dataset Composition}: The dataset contains 870 survey responses collected from 290 distinct students, with each respondent evaluating all three sports (Soccer, Baseball, and Basketball). In this dataset, there are 14 columns representing the survey questions and responses.
    \item \textbf{Class Balance}: The target labels are perfectly balanced (290 observations per sport overall; exactly 203 per sport in the 70\% training split). No class rebalancing or oversampling is needed.
\end{itemize}

\subsection*{Data Splitting}
\begin{itemize}[leftmargin=1.5em]
    \item \textbf{Respondent Isolation (\texttt{id})}: The respondent identifier (\texttt{id}) was used to perform grouped partitioning. All responses from a given student were kept within the same dataset split to prevent data leakage arising from respondent-specific vocabulary, preferences, or writing styles.
    \item \textbf{Train-Test-Validation Split}: The 290 respondent IDs are partitioned into:
    \begin{itemize}
        \item \textbf{Training split (70\%)}: 203 respondents (609 observations) used for exploratory analysis, distribution inspection, and feature engineering.
        \item \textbf{Validation split (15\%)}: 43 respondents (129 observations) reserved strictly for hyperparameter tuning.
        \item \textbf{Internal test split (15\%)}: 44 respondents (132 observations) held out for final evaluation of the best model.
    \end{itemize}

    Following partitioning, the validation and test sets were excluded from all exploratory analysis. Feature engineering, imputation parameters, standardization statistics, and text vocabularies were learned exclusively from the training set and applied to the validation and test sets during evaluation and prediction.
\end{itemize}

\subsection*{Features Types}
\begin{table}[H] \centering \caption{Feature categories in the dataset.} \label{tab:feature_types} \footnotesize \begin{tabular}{p{3cm} p{13cm}} \toprule \textbf{Feature Category} & \textbf{Variables} \\ \midrule Numerical & difficulty\_rating \\ Likert-scale & strategy\_importance, injury\_risk\_comfort, individual\_brilliance \\ Multi-select categorical & play\_location, play\_season, associated\_music \\ Free-text & associated\_foods\_drinks, associated\_color, associated\_song, emotional\_response \\ Semi-structured text & acceptable\_price \\ Target & label \\ \bottomrule \end{tabular} \end{table}


\subsection*{Features Distributions}
The distribution of "ease of starting" ratings appears to be skewed to high values for soccer and basketball and low values for baseball. The Likert-scale features are mostly unimodal, with a slight skew towards higher values for strategy importance and injury risk comfort. Multi-select categorical features show a wide variety of responses, with some locations (e.g. outdoor spaces), music genres (e.g. hip-hop or rap), seasons (e.g. summer) being more popular than others, namely winter and classical/jazz. Free-text features exhibit high variability in length and content, with some respondents providing detailed descriptions while others leave fields blank. Ticket price willingness exhibits some skewness with a few respondents reporting high willingness to pay.


\subsection*{Missing Values, Outliers, and Inconsistencies}
Blank responses are rare overall (\textasciitilde 1.70\%), with the song feature exhibiting the highest vacancy (\textasciitilde 90 entries missing). Rather than discarding sparse text rows, a binary indicator (\texttt{is\_missing}) is appended to each text feature, allowing models to leverage deliberate omission as a predictive signal (e.g., respondents unable to name a song for less mainstream sports). Imputation is also used to handle missing values in numerical and Likert-scale features, with mean imputation for difficulty rating and neutral-value imputation for Likert features. 

Outliers are present in the acceptable price feature, with some respondents reporting extreme willingness to pay; however, no outliers were removed. A $\log(1+x)$ transformation is applied to reduce skewness and mitigate outlier influence. Inconsistencies are also present in this feature with some respondents reporting prices with dollar signs, commas, etc. These are stripped retaining only the numeric value during preprocessing. 


\subsection*{Feature Preprocessing}
\begin{table}[H] \centering \caption{Feature preprocessing summary.} \label{tab:preprocessing} \footnotesize \begin{tabular}{p{2.5cm} p{7.3cm} p{6.5cm}} \toprule \textbf{Feature Group} & \textbf{Preprocessing} & \textbf{Rationale} \\ \midrule Respondent ID & Excluded from predictors; used for grouped partitioning & Prevents overfitting (IDs are not generalizable). \\ Difficulty Rating & Mean imputation ($\mu=5.48$); z-score standardization & Removes scale disparity and handles missing values. \\ Likert Features & Monotonic mapping (1--5); neutral imputation (3) & Preserves ordinal structure while avoiding unnecessary dimensionality. \\ Multi-select Features & Lowercasing; delimiter splitting; multi-hot encoding & Represents multiple simultaneous selections without imposing ordering. \\ Text Features & Lowercasing; punctuation removal; TF--IDF vectorization & Downweights common words and emphasizes discriminative sport-specific terms. \\ Acceptable Price & Regex removal of ``\$'' and commas; conversion of ``thousand'' expressions; median imputation; $\log(1+x)$ transformation & Produces numeric values and reduces extreme positive skewness. \\ Target Label & Map Soccer $\rightarrow 0$, Baseball $\rightarrow 1$, Basketball $\rightarrow 2$ & Creates machine-readable class labels. \\ \bottomrule \end{tabular} \end{table}

\section{Methodology}\label{sec:methodology}
\subsection*{Model Families} 
Three model families were evaluated: Softmax Regression, Random Forests, and Naive Bayes.
\begin{itemize} 
    \item \textbf{Softmax Regression}: Softmax regression was selected as a candidate classification algorithm since response variable contains three mutually exclusive classes: Soccer, Baseball, and Basketball. Softmax regression naturally handles multi-class classification problems by providing a probability distribution over the classes and is a simple, interpretable model. 
    \item \textbf{Naive Bayes}: Naive Bayes was considered for the text-based features which were transformed with TF-IDF as it is effective at handling high-dimensional data. Furthermore, since the model assumes that features are conditionally independent given the target class it can efficiently combine many weak signals from free-text responses such as \texttt{associated\_song}, \texttt{associated\_foods\_drinks}, and \texttt{emotional\_response}. 
    \item \textbf{Random Forest}: Random forest was also selected as a candidate classification as it can handle complex non-linear relationships in multi-class classification. Random forests can therefore capture patterns that softmax regression might miss. Furthermore, by selecting a random subset of features and by performing bagging, random forests reduce the variance in their prediction compared to a single decision tree.
\end{itemize} 

\subsection*{Optimizer, Training Setup, Regularization, and Early Stopping} 
The following section outlines the training setup for each model family.
\begin{itemize} 
    \item \textbf{Softmax Regression}: Softmax regression was implemented using the \texttt{LogisticRegression} classifier from Python's scikit-learn library. The model was trained using the training dataset and the L-BFGS optimization algorithm was used to estimate the model weights. The L-BFGS algorithm is a quasi-Newton method that approximates the Hessian matrix to find the optimal parameters. The maximum number of iterations was set to 1,000 if the algorithm fails to converge. L2 regularization was applied using the default regularization strength (C = 1.0) to penalize large weights and reduce the risk of overfitting.
    \item \textbf{Naive Bayes}: The Naive Bayes models were trained using scikit-learn's \texttt{MultinomialNB} class which learns class prior probabilities and feature likelihoods from the training data. The model was initially trained with the default smoothing parameter $\alpha = 1.0$ to handle zero probabilities for unseen features in the training data. The predictions from the Naive Bayes model on the TF-IDF features were combined with the predictions from a Random Forest model on all features using a weighted average ensemble method. The weights for the ensemble were initially set to 0.5 each and were tuned using the validation set. This ensemble method was used since Naive Bayes is effective at handling high-dimensional text features while Random Forests can capture complex interactions between structured features.
    \item \textbf{Random Forest}: Random forest was implemeneted using the \texttt{RandomForestClassifier} class from Python's scikit-learn library and the \texttt{gini} criterion was used to measure the quality of the splits. The model was trained using the training dataset and the number of trees was set to a default of 100 and later tuned. 
\end{itemize}

\subsection*{Train/Validation Setup} 

As stated in the data exploration section, the data was split using a 70\%/15\%/15\% train / validation / test split based on the unique respondent ID. By splitting based on ID, each respondent appears in only one of the sets which prevents data leakage as responses from the same respondent tend to be correlated. The random seed used for this split was set to 42 to ensure reproducibility.

\subsection*{Hyperparameters}

The following describes the hyperparameters that were tuned for each model family. To tune the hyperparameters, a manual grid search was performed over a range of values for each hyperparameter and the best configuration was selected based on the highest validation accuracy score.
\subsubsection*{Softmax Regression} The hyperparameter that was tuned for the softmax regression models was $C = \frac{1}{\lambda}$ where $\lambda$ is the regularization strength. $C$ was tuned over the values 0.01, 0.1, 0.5, 1.0, 5.0, and 10.0. The default $C$ value used by scikit-learn is 1.0 so the values are chosen to explore a broad range around this value. Based on Table 3 below, the model with the highest validation accuracy was 0.5, 1.0, and 5.0. The model with validation accuracy 0.5 as it has the same validation accuracy with a lower training accuracy which suggests less overfitting.
\begin{table}[h]
    \centering
    \caption{Softmax Regression Hyperparameter Tuning Results}
    \label{tab:softmax_tuning}
    \begin{tabular}{c c c}
        \hline
        \textbf{C Value} & \textbf{Training Accuracy} & \textbf{Validation Accuracy} \\
        \hline
        0.01 & 0.8752 & 0.8837 \\
        0.10 & 0.9097 & 0.9380 \\
        $\boxed{0.50}$ & $\boxed{0.9458}$ & $\boxed{0.9457}$ \\
        1.00 & 0.9589 & 0.9457 \\
        5.00 & 0.9787 & 0.9457 \\
        10.00 & 0.9852 & 0.9380 \\
        \hline
    \end{tabular}
\end{table}
\subsubsection*{Random Forest} For the random forest model, the hyperparameter that was tuned was the maximum depth of the trees over a range of 5, 10, 14, 15, 16, 20, and None. The default value is None but to investigate whether the model is overfitting, we compare the accuracy over the other values. Because bagging is used, the number of trees was not tuned as a hyperparameter as the default value 100 captures enough variance in the data and is less computationally expensive. In table 4 below, the validation accuracy is the highest when the maximum depth is 15. When the maximum depth is less, the model underfits as seen by the low training and validation accuracy, and when it is greater, the model begins to overfit as seen by the high training accuracy and lower validation accuracy.
\begin{table}[h]
    \centering
    \caption{Random Forest Hyperparameter Tuning Results}
    \label{tab:rf_tuning}
    \begin{tabular}{c c c}
        \hline
        \textbf{Maximum Depth} & \textbf{Training Accuracy} & \textbf{Validation Accuracy} \\
        \hline
        5    & 0.9179 & 0.9225 \\
        10   & 0.9754 & 0.9380 \\
        14   & 0.9934 & 0.9380 \\
        $\boxed{15}$ & $\boxed{0.9951}$ & $\boxed{0.9535}$ \\
        16   & 0.9951 & 0.9225 \\
        20   & 0.9951 & 0.9457 \\
        None & 0.9951 & 0.9380 \\
        \hline
    \end{tabular}
\end{table}
\subsubsection*{Naive Bayes} For the Naive Bayes model, the smoothing parameter $\alpha$ was tuned over the values 0.1, 0.5, 1.0, 2.0. The default value is 1.0 so a range of values around 1.0 were investigated. The validation accuracy is highest when $\alpha$ is 0.1 or 0.5 as seen in Table 5 below. For the best model, $\alpha$ was set to 0.1.
\begin{table}[h]
    \centering
    \caption{Multinomial Naive Bayes Hyperparameter Tuning Results}
    \label{tab:nb_tuning}
    \begin{tabular}{c c c}
        \hline
        \textbf{Alpha Value} & \textbf{Training Accuracy} & \textbf{Validation Accuracy} \\
        \hline
        $\boxed{0.1}$ & $\boxed{0.9015}$ & $\boxed{0.9147}$ \\
        $\boxed{0.5}$ & $\boxed{0.9015}$ & $\boxed{0.9147}$ \\
        1.0 & 0.8998 & 0.9070 \\
        2.0 & 0.8949 & 0.9070 \\
        \hline
    \end{tabular}
\end{table}

Finally, the Naive-Bayes values were combined with the Random Forest predictions using a weighted average ensemble method. The weights for the ensemble were tuned over the values 0.5, 0.6, 0.7, 0.8, and 0.9 for the Random Forest model. Table 6 below shows that the most optimal weights were 0.9 and 0.1. 

\begin{table}[h]
    \centering
    \caption{Random Forest and Naive Bayes Ensemble Weight Tuning Results}
    \label{tab:ensemble_tuning}
    \begin{tabular}{c c c}
        \hline
        \textbf{RF Weight} & \textbf{NB Weight} & \textbf{Validation Accuracy} \\
        \hline
        $\boxed{0.9}$ & $\boxed{0.1}$ & $\boxed{0.9457}$ \\
        0.8 & 0.2 & 0.9380 \\
        0.7 & 0.3 & 0.9380 \\
        0.6 & 0.4 & 0.9225 \\
        0.5 & 0.5 & 0.9147 \\
        \hline
    \end{tabular}
\end{table}

\subsection*{Metrics}
The following describes the metrics used to compare the final models from each family. These metrics are reported as one average over all classes for each model.
\begin{itemize}
    \item \textbf{Accuracy}: Accuracy was used as a baseline metric and it is suited for this dataset because the target classes are balanced so it provides a good estimate of the percentage of correct predictions.
    \item \textbf{Precision}: Precision was used as it indicates how often the model is correct when it predicts the class. Using this metric, we can assess whether a model has a high false positive rate.
    \item \textbf{Recall}: Recall was used as it identifies how many of the class's training examples it was able to predict correctly which indicates whether a model tends to miss predictions in a particular class.
    \item \textbf{F1-Score}: F1-Score is the harmonic mean of precision and recall which provides an aggregate of these two metrics.
\end{itemize}

% This section of your report should be at most 3 pages.
% The goal of the Methodology step is to describe how you trained, tuned, and
% evaluated your models in enough detail that a reader could reproduce your results.
% Present the three model families you used and explain why each one suits the
% dataset, drawing on what you found during data exploration. Do not jump straight to
% a deep network, because a simpler model may perform just as well and is easier to
% interpret.
% Make sure to describe how you trained each model so that your results can be
% reproduced. State the optimizer (e.g., SGD or Adam), the learning rate schedule if
% you used one, and any regularization or early stopping. Briefly name the libraries
% you used and describe any custom code you wrote.
% Explain how you split the data into training and validation sets, or how you
% applied cross-validation, so that the reader can trust your comparisons. Tune
% hyperparameters on a validation set that is separate from the training data. Never
% use the test set for training, feature engineering, or tuning.
% Summarize your hyperparameter tuning by stating which hyperparameters you varied,
% how you chose the values to try, and which metric you used to compare them. An
% exhaustive search is not required, but you should show that your choices are
% reasonable.
% Be sure to tune every model family before you compare them. Choosing the best model
% using default hyperparameters and then tuning only that model is unfair to the
% other models, which might improve just as much with tuning. This was the most
% common mistake we saw when grading reports in previous semesters.
% Describe the evaluation metrics you used and justify each one. Report at least one
% metric beyond accuracy, such as precision, recall, or F1, because accuracy alone
% can hide poor performance on one of the classes. Make sure every table and plot has
% appropriate labels and an explanation of what it shows.
% Optionally, you may describe any additional techniques you tried, such as feature
% engineering, clustering, or PCA.
% Appendix~\ref{app:methodology} gives a list of guiding questions to work through
% before you write this section.

% ===========================================
% 5. Results
% ===========================================
\section{Results}\label{sec:results}

% This section should be at most 1 page.
% The goal of the Results step is to show how well your models performed, explain
% which model you chose, and predict how it will do on the unseen test set.
% Report the results of every model using the evaluation metrics you described in the
% Methodology section. Present them in tables or plots that make the comparison easy
% to read.
% Compare performance across the three model families, and state clearly which model
% you chose as your final model. This must be the model implemented in your
% prediction script.
% Analyze the errors your models make. Describe the common misclassifications or
% weaknesses you observed, and include a confusion matrix or a few illustrative
% examples if they help.
% Estimate your final model's accuracy on the unseen test set as a single number. A
% range will earn no marks. Make sure to support your estimate with empirical
% evidence, such as validation results, the stability of your cross-validation
% scores, learning curves, or consistency across random seeds or folds.

\subsection*{Model Performance}

The table below shows the comparison of all three model performances and based on the results, the random forest model performs the best on the training and validation data. Specifically, it shows the highest training accuracy of 0.9951 and a lower accuracy of 0.9535; however, there could be some degree of overfitting as the training accuracy is higher than the validation accuracy. The Random Forest + Naive Bayes model has the second best metric values; however, its high training accuracy is likely indication for overfitting. Softmax regression performed relatively well across all metrics with no clear signs of overfitting. Finally, the Naive Bayes-only model performed poorly and had significantly lower metrics than the other three models indicating underfitting. 

Because the data contains a mixture of feature types and features likely have complex interactions, it is expected that random forest performs well. The Naive Bayes + Random Forest likely performed slightly worse because Naive Bayes on its own performs worse than Random Forest so by taking the weighted average, the Naive Bayes is diluting the Random Forest's good predictions. Finally, softmax regression performed worse than random forest likely because the relationship between the features and target may involve nonlinear patterns.

The metrics for validation across models are all quite high but we suspect it is due to a validation split that is particularly easy to predict and could change if the validation split was different due to the small dataset size. Data leakage is not suspected as we shuffled the report labels in the training set and trained the same models after which the validation accuracy dropped to 0.33.
\begin{table}[h]
    \centering
    \caption{Comparison of Model Performance}
    \label{tab:model_comparison}
    \resizebox{\textwidth}{!}{%
    \begin{tabular}{l c c c c c}
        \hline
        \textbf{Model} & \textbf{Train Acc.} & \textbf{Val. Acc.} & \textbf{Macro Precision} & \textbf{Macro Recall} & \textbf{Macro F1} \\
        \hline
        $\boxed{\text{Random Forest (depth=15)}}$ &
        $\boxed{0.9951}$ &
        $\boxed{0.9535}$ &
        $\boxed{0.9540}$ &
        $\boxed{0.9535}$ &
        $\boxed{0.9537}$ \\
        Softmax Regression (C=10.0) & 0.9458 & 0.9457 & 0.9457 & 0.9457 & 0.9457 \\
        RF + NB Ensemble (90\%/10\%) & 0.9934 & 0.9457 & 0.9464 & 0.9457 & 0.9460 \\
        Multinomial Naive Bayes (alpha=0.1) & 0.9015 & 0.9147 & 0.9157 & 0.9147 & 0.9149 \\
        \hline
    \end{tabular}%
    }
\end{table}

\subsection*{Misclassifications}

In table 8 below, the models struggle distinguishing between classes 0 and 1 which maps to soccer and baseball. Basketball on the other hand is the most easy to distinguish. However, in general, there are only few misclassifications but most appear to be between the categories ``soccer'' and ``baseball''.

\begin{table}[H]
    \centering
    \caption{Confusion Matrices for the Four Classification Models}
    \label{tab:confusion_matrices}
    \small

    % First row
    \begin{minipage}{0.48\textwidth}
        \centering
        \textbf{(a) Softmax Regression}\\[4pt]
        \begin{tabular}{c|ccc}
            & 0 & 1 & 2 \\
            \hline
            0 & 40 & 1 & 2 \\
            1 & 2 & 41 & 0 \\
            2 & 1 & 1 & 41 \\
        \end{tabular}
    \end{minipage}
    \hfill
    \begin{minipage}{0.48\textwidth}
        \centering
        \textbf{(b) Random Forest}\\[4pt]
        \begin{tabular}{c|ccc}
            & 0 & 1 & 2 \\
            \hline
            0 & 40 & 3 & 0 \\
            1 & 2 & 41 & 0 \\
            2 & 1 & 0 & 42 \\
        \end{tabular}
    \end{minipage}

    \vspace{0.6cm}

    % Second row
    \begin{minipage}{0.48\textwidth}
        \centering
        \textbf{(c) Multinomial Naive Bayes}\\[4pt]
        \begin{tabular}{c|ccc}
            & 0 & 1 & 2 \\
            \hline
            0 & 40 & 2 & 1 \\
            1 & 3 & 39 & 1 \\
            2 & 2 & 2 & 39 \\
        \end{tabular}
    \end{minipage}
    \hfill
    \begin{minipage}{0.48\textwidth}
        \centering
        \textbf{(d) RF + NB Ensemble}\\[4pt]
        \begin{tabular}{c|ccc}
            & 0 & 1 & 2 \\
            \hline
            0 & 40 & 3 & 0 \\
            1 & 3 & 40 & 0 \\
            2 & 1 & 0 & 42 \\
        \end{tabular}
    \end{minipage}

    \vspace{0.2cm}
    \footnotesize
    Rows represent actual classes; columns represent predicted classes.
\end{table}

\subsection*{Test Error}


% ================================================
% 8. Contributions and Learning
% ================================================
\section{Contributions and Learning}



This section should contain a list of 3–4 bullets, one for each student. Each
student should write at most three sentences describing their main contributions
and the most important lesson learned.

\begin{itemize}
    \item Sylvia Wong: Throughout this project, I performed part of the exploratory data analysis, programmed the preliminary random forest and softmax models, and wrote the methodology portion of the final report. Throughout this project, the most important lesson that I learned was the importance of exploratory data analysis and careful data cleaning to ensure there was no data leakage. Although model tuning and accuracy is important, when data is cleaned properly, the tuning process becomes a lot easier.
\end{itemize}



% ==================================================
% 9. References
% ==================================================
\section{References}
Use references whenever appropriate (e.g., when you consult textbooks, research
papers, online tutorials, or external code). Cite in the text using \verb|\cite{}|.
For example, let's cite the two references provided in the \texttt{.bib} file
\cite{bishop2006prml,hastie2009esl}.
\bibliographystyle{plain}
\bibliography{references}
% \newpage
% \appendix
% Extra space between the subtitle groups and between the questions in the
% appendices.
\setlist[enumerate,1]{itemsep=1em}
\setlist[enumerate,2]{itemsep=0.6em}
\end{document}
